\subsubsection{Thiết kế mô hình AI}

\paragraph{Lựa chọn Mô hình AI (Model Selection)}

\subparagraph{Xác định Ngữ cảnh và Yêu cầu Nghiệp vụ:}

Để đáp ứng các yêu cầu nghiệp vụ, nhóm đã xem xét hai hướng tiếp cận kỹ thuật:

{\small
\begin{longtable}{|p{3.5cm}|p{5.5cm}|p{5.5cm}|}
\caption{So sánh hai hướng tiếp cận xây dựng mô hình AI}\label{tab:ai-approach-comparison}\\
\hline
\textbf{Tiêu chí} & \textbf{Tự xây dựng Mô hình} & \textbf{Sử dụng Dịch vụ (AaaS)} \\ \hline
\endfirsthead

\multicolumn{3}{c}%
{\tablename\ \thetable\ -- \textit{Tiếp theo trang trước}} \\
\hline
\textbf{Tiêu chí} & \textbf{Tự xây dựng Mô hình} & \textbf{Sử dụng Dịch vụ (AaaS)} \\ \hline
\endhead

\hline \multicolumn{3}{r}{\textit{Tiếp tục trang sau}} \\
\endfoot

\hline
\endlastfoot

Chi phí \& Hạ tầng & Rất cao: Đòi hỏi đầu tư lớn cho hạ tầng GPU/TPU, chi phí R\&D và nhân sự chuyên môn AI. & Thấp/Trung bình: Chi phí vận hành dựa trên API (pay-as-you-go), không yêu cầu chi phí vốn ban đầu cho hạ tầng. \\ \hline

Thời gian Phát triển & Rất dài (Tháng/Năm): Bao gồm thời gian thu thập dữ liệu, gán nhãn, huấn luyện, và tinh chỉnh mô hình. & Nhanh (Tuần/Tháng): Tập trung vào việc tích hợp API, bỏ qua giai đoạn huấn luyện tốn kém. \\ \hline

Yêu cầu Dữ liệu & Cực kỳ lớn: Cần một tập dữ liệu (dataset) khổng lồ, được gán nhãn chất lượng cao về du lịch Việt Nam. & Không yêu cầu: Tận dụng các Mô hình Nền tảng (FMs) đã được huấn luyện trước trên dữ liệu toàn cầu. \\ \hline

Giải quyết Đa phương thức & Cực kỳ phức tạp: Đòi hỏi huấn luyện các mô hình riêng biệt (text, image, video) và xây dựng cơ chế "fusion" (kết hợp) phức tạp. & Có sẵn: Các mô hình AaaS hàng đầu (như Gemini) hỗ trợ đa phương thức (văn bản, hình ảnh, video). \\ \hline

\end{longtable}
}

Phân tích trên cho thấy phương án "Tự xây dựng Mô hình" là không khả thi trong khuôn khổ tài nguyên và thời gian của dự án. Do đó, nhóm lựa chọn hướng tiếp cận AaaS, qua đó tập trung vào trọng tâm của dự án: phát triển kỹ thuật phần mềm và thiết kế kiến trúc tích hợp.

Với hướng tiếp cận AaaS, mô hình AI được chọn phải giải quyết được 4 yêu cầu nghiệp vụ chính của TravelEZ:

\begin{itemize}
    \item \textbf{Tạo Lộ trình Cá nhân hóa (Module 1):} Cần suy luận logic, xử lý nhiều ràng buộc (thời gian, ngân sách, sở thích).
    \item \textbf{Khai phá Xu hướng (Module 6):} Cần xử lý và tổng hợp khối lượng lớn (hàng ngàn) văn bản UGC (reviews, blogs).
    \item \textbf{Phân tích \& Cải thiện Lịch trình (Module 7):} Cần đối chiếu lịch trình có sẵn với dữ liệu xu hướng (từ Module 6).
    \item \textbf{Kiểm duyệt Nội dung (Module 5):} Cần phân tích đa phương thức (văn bản, hình ảnh, và video).
\end{itemize}

\subparagraph{Phân tích Yêu cầu Kỹ thuật và Các Ràng buộc Hệ thống:}

Việc lựa chọn mô hình AI cho TravelEZ được quyết định bởi đặc thù của dữ liệu đầu vào (Input Data) mà hệ thống phải xử lý và yêu cầu về tính đơn giản của kiến trúc (Architectural Simplicity). Dưới đây là 4 thách thức kỹ thuật cốt lõi xác định tiêu chí chọn mô hình:

\textbf{a. Thách thức về Dung lượng Ngữ cảnh:} Đây là rào cản kỹ thuật lớn nhất đối với Module 1 (Gợi ý Lộ trình). Để AI có thể sắp xếp một lịch trình tối ưu, nó không thể hoạt động rời rạc mà phải "nhìn thấy" toàn bộ dữ liệu ứng viên trong một lần suy luận.

Phân tích Dữ liệu đầu vào (Input Data):
\begin{itemize}
    \item Để lập lộ trình cho một thành phố, hệ thống cần truy xuất danh sách khoảng 500 địa điểm (POIs) tiềm năng từ cơ sở dữ liệu.
    \item Mỗi POI chứa thông tin: Tên, Mô tả chi tiết, Giờ mở cửa, Giá vé, Địa chỉ. Trung bình khoảng 700 tokens/POI.
    \item Dữ liệu người dùng (User Input): Sở thích, lịch sử, yêu cầu đặc biệt $\approx$ 2.000 tokens.
    \item System Prompt \& Format instruction: $\approx$ 3.000 tokens.
\end{itemize}

Bài toán tính toán:
\begin{equation}
\text{Tổng Token} = (500 \times 700) + 2.000 + 3.000 \approx 355.000 \text{ (tokens)}
\end{equation}

\textit{Kết luận Tiêu chí 1:} Dung lượng này đã vượt giới hạn của các mô hình tiêu chuẩn (thường là 128k tokens). Để đảm bảo hệ thống hoạt động ổn định khi dữ liệu tăng trưởng (ví dụ: lên 600-800 POIs) và tránh lỗi cắt cụt dữ liệu (truncation error), mô hình bắt buộc phải có Context Window lớn hơn 500.000 tokens.

\textbf{b. Thách thức về Xử lý Đa phương thức:} Đối với Module 5 (Kiểm duyệt nội dung), hệ thống yêu cầu phát hiện các vi phạm (bạo lực, nhạy cảm) trong video do người dùng tải lên. Thách thức nằm ở kiến trúc tích hợp chứ không chỉ là độ thông minh của AI.

Vấn đề Kiến trúc: Nếu chọn mô hình chỉ hỗ trợ xử lý ảnh (Image-only Model), đội dự án buộc phải xây dựng một quy trình Backend phức tạp: Tải video $\rightarrow$ Dùng thư viện (như FFmpeg) cắt video thành hàng trăm ảnh tĩnh (frames) $\rightarrow$ Gửi từng ảnh lên AI $\rightarrow$ Tổng hợp kết quả. Cách làm này gây tốn tài nguyên server, tăng độ trễ và rủi ro lỗi phần mềm.

\textit{Kết luận Tiêu chí 2:} Để đơn giản hóa kiến trúc và giảm tải cho Backend, mô hình được chọn bắt buộc phải hỗ trợ Native Video Input (cho phép gửi trực tiếp file video qua API để mô hình tự xử lý nội bộ). Điều này giúp quy trình kiểm duyệt chỉ còn là 1 API call duy nhất.

\textbf{c. Thách thức về Tổng hợp Dữ liệu Lớn:} Module 6 (Khai phá Xu hướng) yêu cầu hệ thống đọc và tổng hợp thông tin từ hàng ngàn bài review của người dùng (UGC) để đưa ra báo cáo xu hướng.

Phân tích Chi phí \& Hiệu năng: Nếu gọi API riêng lẻ cho từng review (ví dụ: 1.000 reviews = 1.000 API calls), chi phí sẽ rất cao và tốc độ chậm. Giải pháp tối ưu là gom (batch) hàng trăm review vào một prompt duy nhất để AI tổng hợp một lần.

\textit{Kết luận Tiêu chí 3:} Mô hình cần có một phiên bản nhẹ (Lite/Flash) với chi phí thấp nhưng vẫn phải duy trì Context Window lớn (để chứa được batch dữ liệu lớn). Các mô hình giá rẻ thông thường thường bị cắt giảm Context Window xuống thấp (4k-8k tokens), không đáp ứng được yêu cầu nạp nhiều review cùng lúc của TravelEZ.

\textbf{d. Thách thức về Tối ưu hóa Chi phí Vận hành:} Bên cạnh các rào cản về kỹ thuật, chi phí vận hành là yếu tố sống còn đối với khả năng duy trì của dự án, đặc biệt khi hệ thống phải xử lý lượng dữ liệu ngữ cảnh khổng lồ.

Phân tích bài toán: Như đã tính toán, một tác vụ lập lộ trình đơn lẻ tiêu tốn khoảng 355.000 tokens. Đối với Module 6 (Phân tích xu hướng), việc xử lý hàng ngàn bài review cùng lúc cũng yêu cầu lượng token khổng lồ.

Rủi ro: Nếu đơn giá cho mỗi 1 triệu token (Input Price per 1M tokens) quá cao, chi phí vận hành sẽ tăng theo cấp số nhân khi lượng người dùng tăng trưởng, dẫn đến việc dự án không thể duy trì hoặc giá thành dịch vụ quá cao so với chi trả của khách hàng.

\textit{Kết luận Tiêu chí 4:} Hệ thống yêu cầu mô hình được chọn phải có chi phí hợp lý và hiệu quả. Cụ thể, cần ưu tiên các dòng mô hình có phân khúc giá thấp (Low-cost tier) cho các tác vụ xử lý hàng loạt (Batch processing) và tác vụ yêu cầu ngữ cảnh dài, đảm bảo cân bằng giữa hiệu năng xử lý và ngân sách vận hành.

\subparagraph{Khảo sát và Lựa chọn Công nghệ:}

Dựa trên các tiêu chí kỹ thuật đã xác lập, nhóm thực hiện đối chiếu thông số kỹ thuật của ba dòng mô hình phổ biến hiện nay: GPT-5 (OpenAI), Claude 4.5 Sonnet (Anthropic) và Gemini 2.5 (Google).

{\small
\begin{longtable}{|p{3cm}|p{2.5cm}|p{2.5cm}|p{2.5cm}|p{3cm}|}
\caption{Đối chiếu thông số kỹ thuật các mô hình AI}\label{tab:model-comparison}\\
\hline
\textbf{Tiêu chí} & \textbf{Ngưỡng yêu cầu} & \textbf{GPT-5 / 5 Pro} & \textbf{Claude 4.5 Sonnet} & \textbf{Gemini 2.5 Pro / Flash} \\ \hline
\endfirsthead

\multicolumn{5}{c}%
{\tablename\ \thetable\ -- \textit{Tiếp theo trang trước}} \\
\hline
\textbf{Tiêu chí} & \textbf{Ngưỡng yêu cầu} & \textbf{GPT-5 / 5 Pro} & \textbf{Claude 4.5 Sonnet} & \textbf{Gemini 2.5 Pro / Flash} \\ \hline
\endhead

\hline \multicolumn{5}{r}{\textit{Tiếp tục trang sau}} \\
\endfoot

\hline
\endlastfoot

Context Window & Min: 500.000 tokens & 128.000 tokens & 200.000 tokens & 1.048.576 tokens ($\sim$1M) \\ \hline

Video Input & Yêu cầu: Native & Không (Chỉ Image) & Không (Chỉ Image/Doc) & Có (Native Video/Audio) \\ \hline

Batch/Bulk Processing & Yêu cầu: Context cao + chi phí thấp & Context thấp (128k) & Context trung bình (200k) & Context cao (1M) + Batch API \\ \hline

Chi phí input (mỗi 1000 tokens) & Ưu tiên chi phí hợp lý & \$0.00125 & \$0.003 & \$0.00125 (Pro) / \$0.00015 (Flash) \\ \hline

\end{longtable}
}

\textbf{Đánh giá GPT-5 (Base \& Pro):} Mô hình có giới hạn Context Window là 128.000 tokens. Đối chiếu với yêu cầu tính toán (355.000 tokens), mô hình không đáp ứng tiêu chí này. Rủi ro tràn bộ nhớ ngữ cảnh càng cao khi dữ liệu đầu vào biến động tăng. Do đó, phương án này không đạt tiêu chí số 1 dù có mức giá đầu vào là \$0.00125/1k tokens.

\textbf{Đánh giá Claude 4.5 Sonnet:} Mô hình không đáp ứng tiêu chí Context Window (yêu cầu tối thiểu 500.000 tokens). Bên cạnh đó, mô hình cũng không hỗ trợ Native Video (Tiêu chí số 2). Việc triển khai sẽ cần bổ sung hạ tầng xử lý video riêng, làm tăng độ phức tạp của kiến trúc tổng thể. Bên cạnh đó, đây là phương án có chi phí cao nhất trong danh sách khảo sát. Do đó, phương án này cũng không tối ưu về mặt tích hợp hệ thống.

\textbf{Đánh giá Gemini 2.5 (Pro \& Flash):} Hệ sinh thái Gemini 2.5 được lựa chọn vì giải quyết được bài toán cân bằng giữa "Sức mạnh" và "Ngân sách":

\begin{itemize}
    \item \textbf{Về Kỹ thuật:} Context Window lớn (>1 Triệu tokens) và khả năng xử lý Native Video đáp ứng trọn vẹn yêu cầu nghiệp vụ mà không cần chia nhỏ dữ liệu hay cài đặt thêm tool xử lý ảnh.
    \item \textbf{Về Chi phí:} Phiên bản Gemini 2.5 Flash có mức giá cực thấp, chỉ \$0.00015/1k tokens. Mức giá này thấp hơn khoảng 8 lần so với GPT-5 (\$0.00125) và 20 lần so với Claude 4.5 Sonnet (\$0.003).
\end{itemize}

Lựa chọn tối ưu: Nhóm sẽ sử dụng Gemini 2.5 Flash làm mô hình chủ đạo cho các tác vụ xử lý dữ liệu lớn (Module 6, Module 1) để tối đa hóa hiệu quả chi phí, và chỉ sử dụng Gemini 2.5 Pro cho các tác vụ đòi hỏi suy luận phức tạp.

\textbf{Kết luận.} Dựa trên phân tích kỹ thuật (Technical gap analysis), dự án lựa chọn \textbf{Gemini 2.5 Family} làm nền tảng phát triển với chiến lược phân bổ như sau:

\begin{itemize}
    \item \textbf{Gemini 2.5 Flash:} Sử dụng làm mô hình chủ đạo cho Module 1 (Lập lộ trình) và Module 6 (Phân tích xu hướng). Lý do: Tận dụng tốc độ phản hồi nhanh và chi phí thấp nhất để xử lý cửa sổ ngữ cảnh lớn (Context Caching cho 500 POIs) và phân tích dữ liệu hàng loạt.
    
    \item \textbf{Gemini 2.5 Pro:} Sử dụng cho Module 5 (Kiểm duyệt nội dung). Lý do: Cần khả năng suy luận đa phương thức sâu (Deep Multimodal Reasoning) để phát hiện các vi phạm tinh vi trong Video và Hình ảnh mà bản Flash có thể bỏ sót.
\end{itemize}

\paragraph{Lựa chọn Kiến trúc Prompt (Prompt Architecture)}

\subparagraph{Định nghĩa bài toán thiết kế trong bối cảnh AaaS:}

Trong bối cảnh phát triển phần mềm tích hợp LLM, Prompt Engineering đã chuyển dịch từ các thực hành theo kinh nghiệm sang các phương pháp hệ thống hóa. Để đảm bảo độ tin cậy của hệ thống TravelEZ, việc thiết kế Prompt không chỉ đơn thuần là viết câu lệnh, mà là bài toán kiểm soát các Mô thức Lỗi (Failure Modes).

\subparagraph{Xây dựng tiêu chí đánh giá kỹ thuật Prompt:}

Dựa trên khung phân loại lỗi do Tian et al. (2025) đề xuất và được tổng hợp trong báo cáo khảo sát tháng 11/2025 \cite{tian2025llm}, nhóm xác định 3 thách thức kỹ thuật lớn nhất tương ứng với 3 tiêu chí đánh giá:

\begin{enumerate}[label=\alph*., leftmargin=*]
    \item \textbf{Thách thức về đặc tả và ý định:} Các yêu cầu lập lịch trình du lịch thường phức tạp về logic không gian và thời gian. Nếu prompt không rõ ràng, mô hình dễ mắc lỗi "Chỉ dẫn mơ hồ" (Ambiguous Instructions) hoặc "Ràng buộc thiếu đặc tả" (Underspecified Constraints).
    
    \textit{Tiêu chí 1:} Kỹ thuật được chọn phải có khả năng diễn giải chính xác các ý định phức tạp và tuân thủ chặt chẽ logic nghiệp vụ.
    
    \item \textbf{Thách thức về độ tin cậy nội dung:} Mô hình ngôn ngữ có xu hướng sinh ra "Ảo giác" (Hallucination) hoặc sử dụng thông tin sai lệch khi không được kiểm soát. Đối với ứng dụng du lịch, việc bịa đặt địa điểm không tồn tại là lỗi nghiêm trọng thuộc nhóm "Input \& Content Defects".
    
    \textit{Tiêu chí 2:} Kỹ thuật phải tối đa hóa độ trung thực (Fidelity), đảm bảo đầu ra bám sát dữ liệu POIs được cung cấp.
    
    \item \textbf{Thách thức về hiệu suất vận hành:} Các kỹ thuật prompt quá dài hoặc cơ chế suy luận phức tạp có thể gây ra lỗi "Prompt quá khổ" (Excessive Prompt Length) hoặc "Đầu ra không giới hạn" (Unbounded Output), dẫn đến độ trễ cao và chi phí vận hành lớn.
    
    \textit{Tiêu chí 3:} Kỹ thuật phải tối ưu hóa số lượng token (đặc biệt là token đầu ra đắt đỏ) để đảm bảo tính kinh tế.
\end{enumerate}

\subparagraph{Khảo sát các kỹ thuật Prompting:}

Dựa trên tài liệu khảo sát, nhóm thực hiện xem xét 4 kỹ thuật phổ biến nhất hiện nay để tìm ra ứng viên phù hợp:

\begin{itemize}
    \item \textbf{Zero-shot Prompting:} Ra lệnh trực tiếp cho mô hình thực hiện tác vụ mà không cung cấp ví dụ mẫu. Kỹ thuật này dựa hoàn toàn vào tri thức huấn luyện trước của mô hình.
    
    \item \textbf{Few-shot Prompting:} Cung cấp một tập hợp các ví dụ (Input-Output pairs) trong ngữ cảnh để định hướng mô hình (In-context learning).
    
    \item \textbf{Chain-of-Thought (CoT) Prompting:} Yêu cầu mô hình hiển thị quá trình suy luận từng bước (step-by-step reasoning) trước khi đưa ra câu trả lời cuối cùng, giúp giải quyết các bài toán logic phức tạp.
    
    \item \textbf{Structured Prompting:} Sử dụng các khuôn mẫu định sẵn và các ràng buộc về định dạng (như JSON Schema) để hướng dẫn mô hình tuân thủ quy tắc dữ liệu, giảm thiểu lỗi sai cấu trúc.
\end{itemize}

\subparagraph{Đánh giá và luận giải chi tiết:}

Để đảm bảo tính khách quan và nhất quán, nhóm thiết lập thang đo chuẩn hóa (Rubric) gồm 4 mức độ như sau:

\textbf{Tiêu chí 1 - Xử lý Logic:}
\begin{itemize}
    \item \textit{Thấp:} Thường xuyên đưa ra lộ trình bất khả thi hoặc sai lệch hoàn toàn yêu cầu.
    \item \textit{Trung bình:} Hiểu được yêu cầu cơ bản nhưng thất bại khi xử lý các ràng buộc chồng chéo.
    \item \textit{Khá:} Xử lý tốt logic thông thường, đôi khi gặp khó khăn với các ràng buộc phức tạp.
    \item \textit{Cao:} Khả năng suy luận sâu, xử lý tốt các bài toán đa bước (multi-step) về không gian/thời gian.
\end{itemize}

\textbf{Tiêu chí 2 - Độ Tin cậy:}
\begin{itemize}
    \item \textit{Thấp:} Thường xuyên gặp lỗi ảo giác (Hallucination) hoặc sai định dạng JSON.
    \item \textit{Trung bình:} Định dạng đúng nhưng nội dung có thể chứa thông tin chưa chính xác.
    \item \textit{Khá:} Tuân thủ định dạng tốt, giảm thiểu được ảo giác nhờ cơ chế bắt chước hoặc suy luận.
    \item \textit{Cao:} Tuân thủ tuyệt đối định dạng (Schema) và kiểm soát chặt chẽ nguồn dữ liệu đầu vào.
\end{itemize}

\textbf{Tiêu chí 3 - Hiệu suất:}
\begin{itemize}
    \item \textit{Thấp:} Tốn rất nhiều token (đặc biệt là token đầu ra đắt đỏ) và phản hồi chậm.
    \item \textit{Trung bình:} Tiêu tốn tài nguyên ở mức đáng kể (do input dài).
    \item \textit{Khá:} Tối ưu hóa tốt, prompt gọn gàng.
    \item \textit{Cao:} Tiết kiệm nhất, tốc độ phản hồi nhanh nhất.
\end{itemize}

{\small
\begin{table}[htp]
\centering
\caption{Đánh giá các kỹ thuật Prompting theo tiêu chí}\label{tab:prompting-evaluation}
\begin{tabular}{|p{4cm}|p{3.2cm}|p{3.2cm}|p{3.2cm}|}
\hline
\textbf{Kỹ thuật} & \textbf{Tiêu chí 1: Xử lý Logic} & \textbf{Tiêu chí 2: Độ Tin cậy} & \textbf{Tiêu chí 3: Hiệu suất} \\ \hline
Zero-shot & THẤP & THẤP & CAO \\ \hline
Few-shot & TRUNG BÌNH & KHÁ & TRUNG BÌNH \\ \hline
Chain-of-Thought & CAO & KHÁ & THẤP \\ \hline
Structured Prompting & KHÁ & CAO & KHÁ \\ \hline
\end{tabular}
\end{table}
}

\textbf{Luận giải Zero-shot (Thấp - Thấp - Cao):} Đạt mức Cao về hiệu suất vì prompt ngắn nhất, không tốn token ví dụ. Tuy nhiên, logic chỉ ở mức Thấp do thiếu hướng dẫn, dễ dẫn đến lỗi "Ambiguous Instructions".

\textbf{Luận giải Few-shot (Trung bình - Khá - Trung bình):} Đạt mức Khá về độ tin cậy nhờ khả năng bắt chước mẫu (Pattern recognition), giúp giảm sai sót định dạng. Hiệu suất chỉ ở mức Trung bình vì phải nạp thêm nhiều token ví dụ vào đầu vào (Input cost).

\textbf{Luận giải Chain-of-Thought (Cao - Khá - Thấp):} Đạt mức Cao về logic nhờ khả năng suy luận từng bước (Step-by-step), giải quyết tốt các bài toán khó. Tuy nhiên, hiệu suất ở mức Thấp vì mô hình phải sinh ra lượng lớn văn bản suy luận. Với giá token đầu ra của Gemini 2.5 Flash (\$0.0006) đắt gấp 4 lần đầu vào, đây là phương án tốn kém nhất.

\textbf{Luận giải Structured Prompting (Khá - Cao - Khá):} Đạt mức Cao về độ tin cậy nhờ các khuôn mẫu (Templates) và ràng buộc Schema chặt chẽ, loại bỏ lỗi cấu trúc. Logic ở mức Khá, tốt hơn Zero-shot nhưng chưa đạt độ sâu suy luận như CoT. Hiệu suất ở mức Khá vì prompt đi thẳng vào vấn đề, không lan man như CoT và không dài dòng như Few-shot.

\subparagraph{Kết luận và Kiến trúc Đề xuất:}

Dựa trên phân tích đánh đổi, nhóm nhận thấy không có một kỹ thuật đơn lẻ nào đáp ứng trọn vẹn cả 3 tiêu chí. Do đó, nhóm quyết định xây dựng kiến trúc \textbf{Hybrid Structured Prompting} (Lai ghép có cấu trúc).

Kiến trúc này được thiết kế để tận dụng nền tảng vững chắc của Structured Prompting và bổ sung sức mạnh suy luận của CoT theo cách tối ưu nhất:

\textbf{Lớp 1: Suy luận ngầm - Giải quyết tiêu chí 1 \& 3:} Thay vì yêu cầu AI phải xuất ra toàn bộ quá trình suy nghĩ (gây tốn chi phí cho phần Output), nhóm đã lựa chọn sử dụng chỉ thị "Suy nghĩ trong im lặng" hoặc "Lập kế hoạch tổng quát rồi xuất ra dưới dạng JSON". Điều này giúp kích hoạt khả năng suy luận logic của CoT nhưng loại bỏ gánh nặng chi phí của việc sinh văn bản thừa.

\textbf{Lớp 2: Ràng buộc cấu trúc - Giải quyết tiêu chí 2:} Sử dụng định dạng JSON Schema nghiêm ngặt để ngăn chặn lỗi sai định dạng (Undefined Output Format). Kết hợp với các "Luật cấm" (Negative Constraints) để giảm thiểu ảo giác.

\textbf{Lớp 3: Kiểm soát đầu ra:} Để khắc phục tính ngẫu nhiên của mô hình, hệ thống tích hợp bộ Validation ngay sau bước nhận kết quả để phát hiện các lỗi logic hoặc ảo giác còn sót lại.

\textbf{Kết luận:} Kiến trúc Hybrid là lựa chọn tối ưu nhất, cân bằng giữa khả năng xử lý nghiệp vụ phức tạp và hiệu quả kinh tế khi triển khai thực tế.

\paragraph{Thiết kế Prompt Chi tiết theo Module:}

Phần này trình bày thiết kế (design) cho các prompt, áp dụng kiến trúc Structured Prompting đã chọn ở phần trước. Các prompt này được thiết kế như các mẫu (template) động, tuân thủ cấu trúc 3 phần: [Vai trò], [Ngữ cảnh], và [Tác vụ \& Định dạng].

\textit{Lưu ý:} Việc tinh chỉnh chi tiết từng câu chữ (Wording optimization) và đánh giá độ hiệu quả thực tế của prompt sẽ được thực hiện và tối ưu hóa trong giai đoạn phát triển kế tiếp của dự án.

\subparagraph{Thiết kế Prompt cho Module 1: Gợi ý Lộ trình}

\textbf{Lựa chọn Model:} Gemini 2.5 Flash đối với tác vụ lập lịch trình, nhóm quyết định sử dụng phiên bản Gemini 2.5 Flash thay vì phiên bản Pro. Quyết định này dựa trên 2 yếu tố:

\begin{enumerate}[label=\arabic*., leftmargin=*]
    \item \textbf{Tính kinh tế:} Tác vụ này yêu cầu xử lý ngữ cảnh lớn (500 POIs) và trả về định dạng JSON dài. Giá của bản Flash (\$0.00015/1k input) thấp hơn khoảng 8 lần so với bản Pro (\$0.00125/1k), giúp đảm bảo khả năng mở rộng khi lượng người dùng tăng.
    
    \item \textbf{Tốc độ Phản hồi:} Bản Flash có độ trễ thấp hơn, phù hợp với trải nghiệm người dùng cuối (End-user UX) cần nhận kết quả nhanh chóng, trong khi bản Pro thiên về các tác vụ suy luận sâu và chậm hơn.
\end{enumerate}

\textbf{Cơ chế quản lý dữ liệu:} Để giải quyết bài toán ngữ cảnh lớn ($\sim$355.000 tokens) mà không làm tăng chi phí vận hành, nhóm áp dụng tính năng Explicit Context Caching.

\begin{itemize}
    \item \textbf{Cơ chế:} Dữ liệu tĩnh (Danh sách 500 POIs, giờ mở cửa, vị trí) được cache một lần duy nhất trên server của Google.
    \item \textbf{Vận hành:} Mỗi request của người dùng chỉ cần gửi tham chiếu đến cache\_name thay vì gửi lại toàn bộ dữ liệu, giúp giảm chi phí input token xuống mức tối thiểu.
\end{itemize}

\textbf{Cấu trúc Prompt:} Prompt được thiết kế theo kiến trúc Hybrid Structured đã lựa chọn, bao gồm 3 khối thành phần logic:

\begin{itemize}
    \item \textbf{Block 1: Định danh \& Ràng buộc}
    \begin{itemize}
        \item \textit{Vai trò:} Định nghĩa AI là "Chuyên gia Lập kế hoạch Du lịch".
        \item \textit{Luật cấm:} Thiết lập các hàng rào bảo vệ (Guardrails) như: \textit{"Cấm bịa đặt địa điểm (No Hallucination)"}, \textit{"Cấm xếp lịch vào giờ đóng cửa"}.
    \end{itemize}
    
    \item \textbf{Block 2: Tiêm Ngữ cảnh}
    \begin{itemize}
        \item \textit{Static Context:} Trỏ đến Cache ID chứa dữ liệu POIs.
        \item \textit{Dynamic Context:} Chứa thông tin người dùng (User Profile) và các yêu cầu cụ thể (Ngày đi, Sở thích, Ngân sách) được tiêm vào từ Backend tại thời điểm chạy (Runtime).
    \end{itemize}
    
    \item \textbf{Block 3: Chỉ dẫn Tác vụ}
    \begin{itemize}
        \item \textit{Cơ chế Suy luận:} Sử dụng chỉ thị Implicit CoT (ví dụ: \textit{"Think silently about logistics..."}) để yêu cầu mô hình tính toán logic di chuyển ngầm trước khi ra quyết định.
        \item \textit{Định dạng Đầu ra:} Ép buộc trả về dữ liệu theo JSON Schema (được định nghĩa bằng Class trong code) để đảm bảo tính tương thích tuyệt đối với hệ thống phần mềm.
    \end{itemize}
\end{itemize}

\subparagraph{Thiết kế Prompt cho Module 5: Kiểm duyệt Nội dung}

\textbf{Lựa chọn Model:} Gemini 2.5 Pro khác với Module 1 ưu tiên tốc độ và giá, Module 5 ưu tiên độ chính xác và khả năng hiểu đa phương thức (Multimodal Understanding). Nhóm chọn phiên bản Pro vì:

\begin{enumerate}[label=\arabic*., leftmargin=*]
    \item \textbf{Chất lượng Kiểm duyệt:} Việc phát hiện các vi phạm tinh vi trong video (ví dụ: bạo lực ẩn, ký hiệu thù ghét, nội dung nhạy cảm ngữ cảnh) đòi hỏi khả năng suy luận hình ảnh cao cấp mà bản Flash có thể bỏ sót.
    
    \item \textbf{Native Video:} Gemini 2.5 Pro hỗ trợ xử lý native video mượt mà với cửa sổ ngữ cảnh lớn, cho phép phân tích toàn bộ clip mà không cần cắt nhỏ thủ công.
\end{enumerate}

\textbf{Cơ chế Quản lý Dữ liệu.} Sử dụng cơ chế Native Multimodal Input (Đầu vào đa phương thức tự nhiên):

\begin{itemize}
    \item \textbf{Cơ chế:} Thay vì sử dụng các thư viện bên thứ 3 (như FFmpeg) để tách video thành ảnh (gây tốn tài nguyên server), hệ thống gửi trực tiếp File URI (đã upload lên Google File API) vào prompt.
    \item \textbf{Đồng bộ:} Dữ liệu hình ảnh/video được gửi kèm với tiêu đề/mô tả (Caption) của người dùng để AI đánh giá ngữ cảnh tổng thể.
\end{itemize}

\textbf{Cấu trúc Prompt.} Prompt được thiết kế theo kiến trúc Structured Classification (Phân loại có cấu trúc), tập trung vào việc đối chiếu nội dung với bộ quy tắc an toàn (Safety Policy):

\begin{itemize}
    \item \textbf{Block 1: Định danh \& Chính sách}
    \begin{itemize}
        \item \textit{Vai trò:} Định nghĩa AI là "AI Trust \& Safety Officer" (Cán bộ An toàn \& Tin cậy).
        \item \textit{Định nghĩa Chính sách:} Liệt kê rõ ràng các danh mục cấm như:
        \begin{itemize}
            \item Nội dung tình dục/Khỏa thân.
            \item Ngôn từ thù ghét.
            \item Bạo lực/Máu me.
            \item Quảng cáo rác/Lừa đảo.
        \end{itemize}
    \end{itemize}
    
    \item \textbf{Block 2: Tiêm Ngữ cảnh Đa phương thức}
    \begin{itemize}
        \item \textit{Media Input:} [Video File / Image File].
        \item \textit{Text Input:} Văn bản người dùng nhập kèm (Caption, Comments) để phát hiện vi phạm trong ngôn ngữ.
    \end{itemize}
    
    \item \textbf{Block 3: Chỉ dẫn Phân loại \& Định dạng}
    \begin{itemize}
        \item \textit{Cơ chế Phân tích:} Yêu cầu mô hình quét toàn bộ nội dung media và text.
        \item \textit{Định dạng Đầu ra:} Ép buộc trả về JSON với các trường quan trọng để Backend xử lý tự động:
        \begin{itemize}
            \item \texttt{decision}: \texttt{"APPROVED"} / \texttt{"FLAGGED"} / \texttt{"REJECTED"}.
            \item \texttt{violation\_category}: Danh mục vi phạm cụ thể (nếu có).
            \item \texttt{confidence\_score}: Độ tin cậy của phán đoán (0.0 - 1.0). Nếu điểm thấp (< 0.7), hệ thống sẽ chuyển cho người duyệt (Human-in-the-loop).
        \end{itemize}
    \end{itemize}
\end{itemize}

\subparagraph{Thiết kế Prompt cho Module 6: Khai phá Xu hướng}

\textbf{Lựa chọn Model: Gemini 2.5 Flash.} Đối với tác vụ đọc hiểu và tổng hợp thông tin từ văn bản (Text Summarization \& Sentiment Analysis), nhóm ưu tiên sử dụng Gemini 2.5 Flash vì:

\begin{enumerate}[label=\arabic*., leftmargin=*]
    \item \textbf{Tối ưu Chi phí cho Dữ liệu Lớn:} Để phân tích xu hướng, hệ thống cần xử lý hàng nghìn bài đánh giá (Reviews). Với mức giá input cực thấp (\$0.00015/1k tokens), Flash là lựa chọn kinh tế nhất để "đọc" khối lượng văn bản khổng lồ này mà không làm bùng nổ ngân sách.
    
    \item \textbf{Context Window Lớn:} Khả năng xử lý cửa sổ ngữ cảnh hơn 1 triệu tokens cho phép nạp hàng trăm bài review vào một lần gọi (Prompt) duy nhất, giúp AI có cái nhìn tổng quan toàn diện thay vì phân tích rời rạc.
\end{enumerate}

\textbf{Cơ chế quản lý dữ liệu.} Áp dụng kỹ thuật Batch Processing (Xử lý theo Lô) để tối ưu hóa hiệu suất:

\begin{itemize}
    \item \textbf{Cơ chế Gom Lô:} Thay vì gọi API cho từng review riêng lẻ (kém hiệu quả), Backend sẽ gom nhóm khoảng 50-100 reviews thành một "văn bản lớn" (Corpus Chunk) và gửi trong một request duy nhất.
    \item \textbf{Tổng hợp:} Kết quả phân tích từ các lô sẽ được backend tổng hợp lại để ra báo cáo xu hướng cuối cùng.
\end{itemize}

\textbf{Cấu trúc Prompt.} Prompt được thiết kế theo kiến trúc Structured Information Extraction (Trích xuất thông tin có cấu trúc), tập trung vào kỹ thuật Phân tích Cảm xúc theo Khía cạnh (Aspect-Based Sentiment Analysis - ABSA):

\begin{itemize}
    \item \textbf{Block 1: Định danh}
    \begin{itemize}
        \item \textit{Vai trò:} Định nghĩa AI là "Chuyên gia Phân tích Dữ liệu Du lịch" (Travel Data Analyst).
        \item \textit{Mục tiêu:} Phát hiện các mẫu hình lặp lại trong ý kiến khách hàng và xác định cảm xúc chủ đạo.
    \end{itemize}
    
    \item \textbf{Block 2: Tiêm Ngữ cảnh Hàng loạt}
    \begin{itemize}
        \item \textit{Input Data:} Chứa danh sách văn bản thô của 50-100 bài review.
    \end{itemize}
    
    \item \textbf{Block 3: Chỉ dẫn Phân tích \& Định dạng}
    \begin{itemize}
        \item \textit{Cơ chế Phân tích:} Yêu cầu mô hình thực hiện các tác vụ cụ thể:
        \begin{enumerate}[label=\arabic*., leftmargin=*]
            \item Trích xuất từ khóa nổi bật (Trending Keywords).
            \item Phân loại cảm xúc theo khía cạnh (Ví dụ: Dịch vụ, Giá cả, Vệ sinh).
            \item Tóm tắt các điểm đau chính.
        \end{enumerate}
        \item \textit{Định dạng Đầu ra:} Ép buộc trả về JSON để Backend dễ dàng thống kê (Visualize) lên Dashboard.
    \end{itemize}
\end{itemize}

\subparagraph{Thiết kế Prompt cho Module 7: Gợi ý Cải thiện Lịch trình}

\textbf{Lựa chọn Model: Gemini 2.5 Flash.} Nhóm tiếp tục chọn Gemini 2.5 Flash cho tác vụ này vì:

\begin{enumerate}[label=\arabic*., leftmargin=*]
    \item \textbf{Tốc độ Phản hồi:} Tính năng này thường được người dùng kích hoạt khi đang xem lịch trình (On-demand). Người dùng cần nhận được gợi ý sửa đổi ngay lập tức, không thể chờ đợi mô hình Pro suy luận chậm.
    
    \item \textbf{Khả năng So sánh:} Flash đủ mạnh để thực hiện tác vụ so sánh đối chiếu (Mapping \& Comparison) giữa "Lịch trình hiện tại" và "Dữ liệu xu hướng" mà không cần đến khả năng sáng tạo cao cấp của bản Pro.
\end{enumerate}

\textbf{Cơ chế quản lý dữ liệu.} Đây là module tổng hợp thông tin, backend sẽ "tiêm" (inject) đồng thời 3 nguồn dữ liệu vào ngữ cảnh để AI đối chiếu:

\begin{itemize}
    \item \textbf{Source 1 - Lịch trình gốc:} JSON lịch trình hiện tại.
    \item \textbf{Source 2 - Dữ liệu Thị trường:} Kết quả JSON từ Module 6. \textit{Ví dụ: Module 6 báo "Quán A đang bị chê phục vụ tệ", AI sẽ dùng thông tin này để cảnh báo.}
    \item \textbf{Source 3 - Điều kiện Ngoại cảnh:} Dữ liệu thời tiết dự báo, mùa du lịch (Mùa mưa/Mùa khô) được lấy từ API thời tiết.
\end{itemize}

\textbf{Cấu trúc Prompt.} Prompt được thiết kế theo kiến trúc Comparative Audit (Kiểm toán so sánh), sử dụng kỹ thuật suy luận ngầm để tìm ra các điểm bất hợp lý (Gaps):

\begin{itemize}
    \item \textbf{Block 1: Định danh \& Mục tiêu}
    \begin{itemize}
        \item \textit{Vai trò:} Định nghĩa AI là "Chuyên gia Tối ưu hóa Trải nghiệm Du lịch" (Experience Optimizer).
        \item \textit{Nhiệm vụ:} Tìm ra các rủi ro (Risk) trong lịch trình và đề xuất phương án thay thế tốt hơn, nhưng \textbf{không được thay đổi cấu trúc cốt lõi} nếu không cần thiết.
    \end{itemize}
    
    \item \textbf{Block 2: Tiêm Ngữ cảnh Đa chiều}
    \begin{description}
        \item[Current Itinerary:] \texttt{\{json\_data\}}
        \item[Trend Warnings:] \texttt{\{bad\_reviews\_list\_from\_module\_6\}}
        \item[Weather Forecast:] \texttt{\{weather\_data\}} (Ví dụ: "Ngày 2: Mưa to")
    \end{description}
    
    \item \textbf{Block 3: Chỉ dẫn Tác vụ \& Định dạng}
    \begin{itemize}
        \item \textit{Cơ chế Suy luận Ngầm:} Chỉ thị \textit{"Cross-check silently"}. Yêu cầu AI so khớp từng hoạt động trong lịch trình với dữ liệu xu hướng và thời tiết.
        \begin{itemize}
            \item \textit{Logic:} Nếu [Hoạt động ngoài trời] AND [Dự báo mưa] $\rightarrow$ Đề xuất đổi sang Bảo tàng/Indoor.
            \item \textit{Logic:} Nếu [Điểm đến A] AND [Trend xấu từ Module 6] $\rightarrow$ Đề xuất điểm B cùng loại (Alternative).
        \end{itemize}
        \item \textit{Định dạng Đầu ra (Schema Enforcement):} Trả về JSON danh sách các đề xuất cải tiến.
    \end{itemize}
\end{itemize}
